\documentclass[a4paper,10pt,french]{article}
\usepackage{../../CLT2425}
\geometry{a4paper,tmargin=1cm,bmargin=0.5cm,lmargin=1cm,rmargin=1cm}
%\renewcommand{\familydefault}{\sfdefault}
% \DeclareSymbolFont{numbers}{T1}{cmss}{m}{n}
%%  \SetSymbolFont{numbers}{}{T1}{phv}{bx}{n}
%  \DeclareMathSymbol{0}\mathalpha{numbers}{"30}
%  \DeclareMathSymbol{1}\mathalpha{numbers}{"31}
%  \DeclareMathSymbol{2}\mathalpha{numbers}{"32}
%  \DeclareMathSymbol{3}\mathalpha{numbers}{"33}
%  \DeclareMathSymbol{4}\mathalpha{numbers}{"34}
%  \DeclareMathSymbol{5}\mathalpha{numbers}{"35}
%  \DeclareMathSymbol{6}\mathalpha{numbers}{"36}
%  \DeclareMathSymbol{7}\mathalpha{numbers}{"37}
%  \DeclareMathSymbol{8}\mathalpha{numbers}{"38}
%  \DeclareMathSymbol{9}\mathalpha{numbers}{"39}

\renewcommand{\baselinestretch}{1.2}


\setlength{\columnsep}{30pt}
\begin{document}
\pagestyle{empty}

\titrebis{0.45}{Feuille 5 - Produit par un scalaire}

\bigskip



\textbf{{\large Exercice 1.} }\\
Exprimer le vecteur  $\Vj{AC}$ en fonction du vecteur $\Vj{AB}$ dans chacun des cas suivant :\\
\\
\begin{tabular}{p{0.5\textwidth} p{0.5\textwidth} }
\textcircled{1} 

\psset{xunit=1.685067139411767cm,yunit=1.0cm,algebraic=true,dotstyle=o,dotsize=3pt 0,linewidth=0.8pt,arrowsize=3pt 2,arrowinset=0.25}
\begin{pspicture*}(-1.71,-0.5)(2.44,0.5)
\psaxes[labelFontSize=\scriptstyle,xAxis=true,yAxis=false,labels=none,Dx=0.5,Dy=0.5,ticksize=-1.5pt 1.5pt]{-}(0,0)(-1.71,-0.5)(2.44,0.5)
\psdots[dotstyle=*](0,0)
\rput[b](0,0.15){A}
\psdots[dotstyle=*](2,0)
\rput[b](2,0.15){B}
\psdots[dotstyle=*](-1.5,0)
\rput[b](-1.5,0.15){C}
\end{pspicture*}

Les vecteurs $\Vj{AC}$ et $\Vj{AB}$ ont même direction. 

Les vecteurs $\Vj{AC}$ et $\Vj{AB}$ ont  \begin{tabular}{l} $\square$ même sens \\ $\square$  sens opposé \end{tabular}.

On a : AC= .... et AB = ....

\medskip

Donc la longueur AC vaut \ $\dfrac{\ ...\ }{\ ... \ }$ \ celle de AB.

\medskip

On a donc :  $\Vj{AC} = .... \ \dfrac{\ ...\ }{\ ...\ } \ \Vj{AB} $
&
\textcircled{2}

\psset{xunit=2.7127027812500035cm,yunit=1.0cm,algebraic=true,dotstyle=o,dotsize=3pt 0,linewidth=0.8pt,arrowsize=3pt 2,arrowinset=0.25}
\begin{pspicture*}(-0.18,-0.5)(2.4,0.5)
\psaxes[labelFontSize=\scriptstyle,xAxis=true,yAxis=false,labels=none,Dx=0.5,Dy=0.5,ticksize=-1.5pt 1.5pt]{-}(0,0)(-0.18,-0.5)(2.4,0.5)
\psdots[dotstyle=*](0,0)
\rput[b](0,0.15){A}
\psdots[dotstyle=*](0.51,0)
\rput[b](0.55,0.15){B}
\psdots[dotstyle=*](2,0)
\rput[b](2,0.15){C}
\end{pspicture*}

Les vecteurs $\Vj{AC}$ et $\Vj{AB}$ ont même direction. 

Les vecteurs $\Vj{AC}$ et $\Vj{AB}$ ont  \begin{tabular}{l} $\square$ même sens \\ $\square$  sens opposé \end{tabular}.

On a : AC= ... et AB = ....

\medskip

Donc la longueur AC \ $\dfrac{\ ...\ }{\ ... \ }$ \ celle de AB.

\medskip

On a donc :  $\Vj{AC} = .... \ \ \dfrac{\ ...\ }{\ ... \ } \   \Vj{AB} = .... \  \Vj{AB}  $


\\
& \\
\textcircled{3}

\psset{xunit=2.765071560810815cm,yunit=1.0cm,algebraic=true,dotstyle=o,dotsize=3pt 0,linewidth=0.8pt,arrowsize=3pt 2,arrowinset=0.25}
\begin{pspicture*}(-1.69,-0.5)(0.84,0.5)
\psaxes[labelFontSize=\scriptstyle,xAxis=true,yAxis=false,labels=none,Dx=0.5,Dy=0.5,ticksize=-1.5pt 1.5pt]{-}(0,0)(-1.69,-0.5)(0.84,0.5)
\psdots[dotstyle=*](0,0)
\rput[b](0,0.15){A}
\psdots[dotstyle=*](0.51,0)
\rput[b](0.55,0.15){B}
\psdots[dotstyle=*](-1.5,0)
\rput[b](-1.5,0.15){C}
\end{pspicture*}

&
\textcircled{4}

\psset{xunit=3.3622231654929626cm,yunit=1.0cm,algebraic=true,dotstyle=o,dotsize=3pt 0,linewidth=0.8pt,arrowsize=3pt 2,arrowinset=0.25}
\begin{pspicture*}(-0.75,-0.5)(1.33,0.5)
\psaxes[labelFontSize=\scriptstyle,xAxis=true,yAxis=false,labels=none,Dx=0.5,Dy=0.5,ticksize=-1.5pt 1.5pt]{-}(0,0)(-0.75,-0.5)(1.33,0.5)
\psdots[dotstyle=*](0,0)
\rput[b](0,0.15){A}
\psdots[dotstyle=*](1,0)
\rput[b](1,0.15){B}
\psdots[dotstyle=*](-0.5,0)
\rput[b](-0.5,0.15){C}
\end{pspicture*}
\\
& \\
 \textcircled{5}
 
 \psset{xunit=1.7091969791169475cm,yunit=1.0cm,algebraic=true,dotstyle=o,dotsize=3pt 0,linewidth=0.8pt,arrowsize=3pt 2,arrowinset=0.25}
\begin{pspicture*}(-0.23,-0.5)(3.86,0.5)
\psaxes[labelFontSize=\scriptstyle,xAxis=true,yAxis=false,labels=none,Dx=0.5,Dy=0.5,ticksize=-1.5pt 1.5pt]{-}(0,0)(-0.23,-0.5)(3.86,0.5)
\psdots[dotstyle=*](0,0)
\rput[b](0,0.15){A}
\psdots[dotstyle=*](3.5,0)
\rput[b](3.5,0.15){B}
\psdots[dotstyle=*](1,0)
\rput[b](1,0.15){C}
\end{pspicture*}
 &
 \textcircled{6}
   
\psset{xunit=1.9674547644230795cm,yunit=1.0cm,algebraic=true,dotstyle=o,dotsize=3pt 0,linewidth=0.8pt,arrowsize=3pt 2,arrowinset=0.25}
\begin{pspicture*}(-0.24,-0.5)(3.31,0.5)
\psaxes[labelFontSize=\scriptstyle,xAxis=true,yAxis=false,labels=none,Dx=0.5,Dy=0.5,ticksize=-1.5pt 1.5pt]{-}(0,0)(-0.24,-0.5)(3.31,0.5)
\psdots[dotstyle=*](0,0)
\rput[b](0,0.15){A}
\psdots[dotstyle=*](3,0)
\rput[b](3,0.15){B}
\psdots[dotstyle=*](0.5,0)
\rput[b](0.5,0.15){C}
\end{pspicture*}
\end{tabular}


\bigskip

\bigskip



\textbf{{\large Exercice 2.}}\\
Dans chacun des cas suivant, repr\'esenter le vecteur $\Vj{AC}$ tel que :\\
\\
\begin{tabular}{p{0.5\textwidth} p{0.5\textwidth} }
 \textcircled{1} $\Vj{AC} = -\dfrac{3}{2}\Vj{AB}$
 
 \psset{xunit=1.535684879929396cm,yunit=1.0cm,algebraic=true,dotstyle=o,dotsize=3pt 0,linewidth=0.8pt,arrowsize=3pt 2,arrowinset=0.25}
\begin{pspicture*}(-1.78,-0.5)(2.78,0.5)
\psaxes[labelFontSize=\scriptstyle,xAxis=true,yAxis=false,labels=none,Dx=0.5,Dy=0.5,ticksize=-1.5pt 1.5pt]{-}(0,0)(-1.78,-0.5)(2.78,0.5)
\psdots[dotstyle=*](0,0)
\rput[bl](0,0.15){A}
\psdots[dotstyle=*](1,0)
\rput[bl](1,0.15){B}
\end{pspicture*}

Les vecteurs $\Vj{AC}$ et $\Vj{AB}$ ont même direction. 

\medskip

Les vecteurs $\Vj{AC}$ et $\Vj{AB}$ ont  \begin{tabular}{l} $\square$ même sens \\ $\square$  sens opposé \end{tabular}.

\medskip

On a : AB= ... , or AC =$\  \dfrac{\ ...\ }{\ ...\ } \quad \times$ AB

\bigskip

Donc la longueur AC vaut \ $\dfrac{\ ...\ }{\ ..\ .}\  \times ... = $ .... 

 &
 \textcircled{2} $\Vj{AC} = \dfrac{7}{8}\Vj{AB}$
 
 \psset{xunit=1.2900832558985085cm,yunit=1.0cm,algebraic=true,dotstyle=o,dotsize=3pt 0,linewidth=0.8pt,arrowsize=3pt 2,arrowinset=0.25}
\begin{pspicture*}(-1.21,-0.5)(4.21,0.5)
\psaxes[labelFontSize=\scriptstyle,xAxis=true,yAxis=false,labels=none,Dx=0.5,Dy=0.5,ticksize=-1.5pt 1.5pt]{-}(0,0)(-1.21,-0.5)(4.21,0.5)
\psdots[dotstyle=*](-0.5,0)
\rput[bl](-0.5,0.15){A}
\psdots[dotstyle=*](3.5,0)
\rput[bl](3.5,0.15){B}
\end{pspicture*}

Les vecteurs $\Vj{AC}$ et $\Vj{AB}$ ont même direction. 

\medskip

Les vecteurs $\Vj{AC}$ et $\Vj{AB}$ ont  \begin{tabular}{l} $\square$ même sens \\ $\square$  sens opposé \end{tabular}.

\medskip

On a : AB= ... , or AC =$\  \dfrac{\ ...\ }{\ ...\ } \quad \times$ AB

\bigskip

Donc la longueur AC vaut \ $\dfrac{\ ...\ }{\ ..\ .}\  \times ... = $ .... 

\\
& \\
 \textcircled{3} $\Vj{AC} = \dfrac{1}{4}\Vj{AB}$
 
\psset{xunit=1.685067139411767cm,yunit=1.0cm,algebraic=true,dotstyle=o,dotsize=3pt 0,linewidth=0.8pt,arrowsize=3pt 2,arrowinset=0.25}
\begin{pspicture*}(-1.71,-0.5)(2.44,0.5)
\psaxes[labelFontSize=\scriptstyle,xAxis=true,yAxis=false,labels=none,Dx=0.5,Dy=0.5,ticksize=-1.5pt 1.5pt]{-}(0,0)(-1.71,-0.5)(2.44,0.5)
\psdots[dotstyle=*](0,0)
\rput[b](0,0.15){A}
\psdots[dotstyle=*](2,0)
\rput[b](2,0.15){B}
%\psdots[dotstyle=*](-1.5,0)
%\rput[b](-1.5,0.15){C}
\end{pspicture*}

 & 
 \textcircled{4} $\Vj{AC} = -\dfrac{1}{2}\Vj{AB}$
 
\psset{xunit=0.6217679467305537cm,yunit=1.0cm,algebraic=true,dotstyle=o,dotsize=3pt 0,linewidth=0.8pt,arrowsize=3pt 2,arrowinset=0.25}
\begin{pspicture*}(-3.68,1.3)(7.58,2.6)
\psaxes[labelFontSize=\scriptstyle,xAxis=true,yAxis=false,labels=none,Dx=1,Dy=1,ticksize=-1.5pt 1.5pt,subticks=2]{-}(0,0)(-3.68,1.3)(7.58,2.6)
\psplot{-3.68}{7.58}{(--6-0*x)/3}
\psdots[dotstyle=*](0,2)
\rput[bl](0.1,2.15){A}
\psdots[dotstyle=*](3,2)
\rput[bl](3.0,2.15){B}
\end{pspicture*}
\\
 \textcircled{5} $\Vj{AC} = \dfrac{2}{3}\Vj{AB}$
 
\psset{xunit=1.685067139411767cm,yunit=1.0cm,algebraic=true,dotstyle=o,dotsize=3pt 0,linewidth=0.8pt,arrowsize=3pt 2,arrowinset=0.25}
\begin{pspicture*}(-1.71,-0.5)(2.44,0.5)
\psaxes[labelFontSize=\scriptstyle,xAxis=true,yAxis=false,labels=none,Dx=0.5,Dy=0.5,ticksize=-1.5pt 1.5pt]{-}(0,0)(-1.71,-0.5)(2.44,0.5)
\psdots[dotstyle=*](0,0)
\rput[b](0,0.15){A}
\psdots[dotstyle=*](1.5,0)
\rput[b](1.5,0.15){B}
%\psdots[dotstyle=*](-1.5,0)
%\rput[b](-1.5,0.15){C}
\end{pspicture*}
&
 \textcircled{6} $\Vj{AC} = -\dfrac{4}{3}\Vj{AB}$
 
\psset{xunit=1cm,yunit=1.0cm,algebraic=true,dotstyle=o,dotsize=3pt 0,linewidth=0.8pt,arrowsize=3pt 2,arrowinset=0.25}
\begin{pspicture*}(-2,1.3)(6,2.6)
\psaxes[labelFontSize=\scriptstyle,xAxis=true,yAxis=false,labels=none,Dx=1,Dy=1,ticksize=-1.5pt 1.5pt,subticks=2]{-}(0,0)(-3.68,1.3)(7.58,2.6)
\psplot{-3.68}{6}{(--6-0*x)/3}
\psdots[dotstyle=*](2,2)
\rput[bl](2.1,2.15){A}
\psdots[dotstyle=*](5,2)
\rput[bl](5.1,2.15){B}
\end{pspicture*}
\end{tabular}

\end{document}